% ============================================
% Johns Hopkins Letter Template
% Assistant Professor Cover Letter – Kutztown University of Pennsylvania
% ============================================

\documentclass[11pt,letterpaper]{letter}

\usepackage{graphicx}
\usepackage{eso-pic}
\usepackage{geometry}
\usepackage{microtype}
\usepackage[hidelinks]{hyperref}

% ----- Page Setup -----
\geometry{left=25mm,right=25mm,top=25mm,bottom=25mm}

% ----- Logo Placement (first page only) -----
\newcommand{\PutJHULogoFG}[1]{%
  \AddToShipoutPictureFG*{%
    \AtPageUpperLeft{%
      \hspace{10mm}%
      \raisebox{-38mm}{%
        \includegraphics[width=6cm]{#1}%
      }%
    }%
  }%
}

% ----- Sender Information -----
\signature{Decory Edwards}
\address{Department of Economics \\ Johns Hopkins University \\ Baltimore, MD 21218 \\ \texttt{dedwar65@jh.edu}}

\begin{document}
\PutJHULogoFG{Images/jhulogo.png}

\begin{letter}{Search Committee \\
Department of Business Administration \\
Kutztown University of Pennsylvania \\
425 Baldy Road \\
Kutztown, PA 19530 \\ USA}

\opening{Dear Members of the Search Committee,}

I am writing to express my interest in the tenure-track Assistant Professor position in Economics in the Department of Business Administration at Kutztown University of Pennsylvania, beginning Fall 2026. I am a Ph.D. candidate in Economics at Johns Hopkins University, advised by Professors Christopher D.~Carroll, Nicholas W.~Papageorge, and Stelios Fourakis, and I expect to receive my degree in May 2026. My research fields are macroeconomics, wealth and income dynamics, and computational economics.

My research agenda focuses on understanding the determinants of wealth inequality and the mechanisms through which heterogeneity in returns to wealth shapes aggregate outcomes. In my job-market paper, I estimate the degree of heterogeneity in individual portfolio returns using micro-data from the Survey of Consumer Finances and simulate a structural model in which households face idiosyncratic return risk. I show that allowing for persistent heterogeneity in returns substantially improves the model’s ability to match the empirical distribution of wealth, offering a tractable way to connect micro-level return dispersion to macroeconomic inequality.

In a second project, I use the Health and Retirement Study to examine the relationship between interpersonal trust and portfolio performance. Building on the literature that finds a hump-shaped relationship between trust and income, I show that a similar relationship holds for individual returns to wealth measured in the HRS data, highlighting a behavioral channel through which social attitudes shape saving and investment outcomes. This work also provides new estimates of the persistent component of returns to net worth, which motivate the calibration of heterogeneity in my structural model.

Kutztown University’s AACSB-accredited Department of Business Administration and its commitment to high-quality undergraduate and graduate instruction make this position an excellent fit. I am particularly drawn to the opportunity to teach \textit{Principles of Microeconomics}, \textit{Principles of Macroeconomics}, \textit{Statistics}, and additional quantitative or analytics-focused courses as needed. The department’s emphasis on teaching effectiveness, professional development, and continued scholarly growth aligns closely with my approach to the teacher–scholar model. I also welcome the opportunity to incorporate artificial intelligence tools into the classroom, consistent with the department’s preferred qualifications and Kutztown’s interest in innovative instructional methods.

\textbf{Teaching.} My teaching experience centers on undergraduate macroeconomics and an upper-level macro policy seminar, where I have led weekly discussion sections and held regular office hours for classes ranging from 16 to 40 students. My approach is student-centered: I aim to help students “convince themselves” that they have moved from confusion to understanding by holding structured office-hour coaching that builds trust and confidence. At Kutztown University, I am prepared to teach \textit{Principles of Microeconomics}, \textit{Principles of Macroeconomics}, \textit{Statistics}, and data-analytics-oriented courses in addition to other quantitative offerings aligned with departmental needs. I am also willing to teach during daytime and evening hours and to integrate AI and data-driven tools into my curriculum when appropriate.

I believe I would be a strong fit because my empirical and computational training allows me to (i) provide rigorous quantitative instruction using real-world datasets, (ii) develop analytics-focused content that supports business and economics students, and (iii) maintain a research agenda suitable for peer-reviewed publication while contributing meaningfully to departmental service expectations. I am also strongly committed to teaching diverse student populations and fostering an inclusive learning environment consistent with Kutztown University’s mission.

I would be honored to join the Department of Business Administration at Kutztown University of Pennsylvania. Thank you for your consideration; I welcome the opportunity to discuss my fit with your department at your convenience.

\closing{Sincerely,}

\end{letter}
\end{document}
